\section{Background and Motivation}

\subsection{PostgreSQL extended statistics}

Ordinary PostgreSQL column statistics summarize one attribute at a time.  A
planner that combines such summaries can therefore miss correlations between
predicates, or approximate them with an independence assumption.  PostgreSQL
extended statistics add multivariate state for a selected group of columns.
The supported fragment in this paper uses two native kinds: most-common-value
(MCV) statistics, which record frequent combinations and their frequencies,
and functional dependencies (``dependencies''), which summarize directional
relationships between columns~\cite{postgresql16docs,postgresql16source}.

The SQL definition and the data-derived state are different objects.
\texttt{CREATE STATISTICS} registers a statistics definition for a relation,
its attributes, and one or more requested kinds.  PostgreSQL records that
definition in \texttt{pg\_statistic\_ext}, including its object identity and
keys.  \texttt{ANALYZE} subsequently derives data-dependent summaries from
the relation and stores the resulting native values in
\texttt{pg\_statistic\_ext\_data}.  The MCV or dependency payload is thus a
property of a definition \emph{under a data realization}, not a constant
property of the column names.

The statistics target controls how much detail PostgreSQL attempts to retain
for the requested statistics.  It affects the resulting payload and its
resource requirements, so it is part of the configuration identity in the
current system.  The product fixes an externally supplied global target
($T=100$ by default) before capture and advice; it does not optimize a
per-candidate target.  A registered definition can also have no requested
native payload after \texttt{ANALYZE}; the advisor represents this legitimate
outcome as \texttt{ABSENT\_NATIVE} rather than fabricating an empty value.

For the rest of the paper, it is useful to name the two kinds of state.  The
\emph{definition state} comprises the relation, the ordered attribute keys,
the requested mechanism, the target and the statistics-object identity.  The
\emph{payload state} comprises the native MCV bytes, dependency bytes, or an
explicit native absence produced for that definition under one realization.
Physical extended-statistics evaluation ordinarily couples both: the
definition must exist in the catalog before PostgreSQL can collect or consume
its payload.

\subsection{Why configuration evaluation is difficult}

The straightforward way to evaluate a candidate design is physical.  For a
design, a tuning process creates the requested statistics definitions, applies
the target with \texttt{ALTER STATISTICS ... SET STATISTICS} when needed, runs
\texttt{ANALYZE}, and executes the workload with native \texttt{EXPLAIN}.
Dropping or replacing definitions is then needed before testing another
configuration.  This path preserves PostgreSQL semantics, but it couples a
configuration search to repeated catalog mutation and data-dependent payload
materialization.  The issue is repeated materialization work, not an
unsupported claim about wall-clock performance.

Extended-statistics utility is also contextual.  With several candidate
objects active, PostgreSQL may change which objects are applicable, their
precedence or winner selection, the clauses consumed by one object, and the
way MCV and dependency information is composed.  The resulting cardinality
estimate can therefore depend on the rest of the active configuration.
Consequently, the usefulness of a candidate cannot generally be treated as a
configuration-independent scalar.  This is an established interaction
phenomenon; the evidence for it in the supported PostgreSQL fragment is
reported later, not asserted as a new theorem here.

\subsection{Existing what-if and statistics-management approaches}

Several prior-art families define the space.  MNSA, StatAdvisor, and DB2
establish workload-driven selection, maintenance concerns, and
statistics-bearing objects~\cite{chaudhuri2001statistics,elhelw2009statadvisor,db2statviews};
the selection problem is not new.  AutoAdmin, SQL Server's Database Engine
Tuning Advisor, and HypoPG establish what-if physical design and backend-local
hypothetical objects~\cite{chaudhuri1998whatif,microsoftdta,hypopg}.

Oracle pending/imported statistics and PostgreSQL 19's
\texttt{pg\_restore\_extended\_stats()} establish statistics testing and
transport~\cite{oraclependingstats,oracleimportstats,postgresql19restore}.
Accordingly, workload selection, what-if analysis, hypothetical objects, and
statistics transport or restoration are not claims of this paper; the detailed
boundaries are given in Section~9.

\subsection{Design alternatives}

These precedents leave a narrower systems capability required by our
PostgreSQL setting.  This paper focuses on a backend-local evaluation
substrate in which a frozen catalog of native MCV/FD candidate states can be
registered, selectively activated in representable configurations, and
consumed by PostgreSQL's native estimator without physically materializing
each design.  The scope is deliberately limited to the supported PostgreSQL
16.14, one-base-relation, fixed-target, arity-two fragment; it is not a claim
about generic what-if optimization or arbitrary PostgreSQL statistics.

There are three conceptual alternatives.  \emph{Physical-per-design}
evaluation keeps native semantics and persistent catalog state, but repeats
definition creation, payload collection, and planner evaluation for each
configuration.  \emph{External estimator emulation} makes hypothetical
configurations easy to represent after one acquisition, but duplicates a
selected subset of PostgreSQL's applicability, precedence, clause
consumption, and MCV/FD composition rules outside the database.  An earlier
CE-Replay project line used this approach for a bounded research fragment; it
is historical context here, not the runtime mechanism of the current advisor.

\emph{Native-state virtualization} acquires one frozen realization, retains
the native candidate states, exposes a transient backend-local membership and
order, and asks PostgreSQL to perform the native estimate.  The key design
choice is therefore to virtualize statistics state, not cardinality-estimator
semantics.  Section~3 formalizes the resulting definition, realization, and
budget contract; Sections~4 and~5 then describe the backend-local substrate
and native evaluation boundary.
